\chapter{Introducció}
Des de ben petit que m'han agradat molt les matemàtiques i quan acabi batxillerat vull fer-ne el grau. Per a qualsevol estdudiant que vulgui cursar una carrera com Matemàtiques, com jo, es troba que la majoria d'universitats ponderen Matemàtiques (MAT), Matemàtiques Aplicades a les Ciències Socials (MACs) i Física. Ara bé, si es cursa el batxillerat cientfic-tècnic, en el qual les matemàtiques és assignatura obligatoria, és molt complicat en la majoria de centres educatius fer MACs alhora. No existeix actualment cap llibre o guia en què s'expliquin els temes que entren a les PAU de les MACs i no s'estudien en les MAT. Per això, em plantejo la necessitat de fer un manual dirigit a preparar els/les estudiants/es del cientfic-tècnic que vulguin fer l'examen de les MACs en les PAU.

En aquest projecte tinc l'interès de:
\begin{enumerate}[·]
 \item Ampliar el meu coneixement de les matemàtiques.
 \item Aprendre a aprendre de manera autònoma.
 \item Aprendre a explicar.
 \item Aprendre a utilitzar eines avançades, com \textbf{Linux}~\cite{linux-foundation}, \textbf{Git}~\cite{github} i \textbf{LaTeX}~\cite{latex-project}, per a treballar de manera més organitzada.
\end{enumerate}

En resum, aquest treball resol una necessitat real per a molts estudiants, també és una oportunitat per aprendre eines avançades i a afrontar un treball de manera professional.
